\scanpage{53}
\mspage{59}

Ainsi, les foncteurs propres $f:X\longrightarrow Y$ se caractérisent de \uline{quatre} façons remarquables

a) $f_*$ commute à tout changement de base ($\forall$ dérivateur \inserted{$\unsure{\mathbb D|\underline D}$} exact à g.)

b) le chang$^{\mathrm{mt}}$ de base $f^*$ commute à tout $g_!$ ($g:\deleted{\unsure{X}}Y'\unsure{\relbar\joinrel\relbar|\longrightarrow}Y$, $\mathbb D$ dér. exact à droite)

c) $\forall\deleted{\unsure{x}}y\in\mathrm{Ob}\,\deleted{\unsure{X}}Y$, $\underset{\text{fibre de }X\text{ en }y}{X_{\deleted{\unsure{x}}y}}\longrightarrow X_{/y\deleted{\unsure{f(x)}}}$ est une hot. équiv. costricte \inserted{i.e. est un foncteur homotop$^{\mathrm t}$ cofinal}

d) $\forall x\in\mathrm{Ob}\,X$, \deleted{\unsure{$x\mapsto$}} le foncteur induit sur les colocalisés

\[
{}_{x\backslash}X_{\deleted{\unsure{/x}}}\longrightarrow{}_{f(x)\backslash}Y_{\deleted{\unsure{/f(x)}}},
\]

est à fibres contractiles.

\[ \rule{1.3cm}{0.4pt} \]

Autre situation \unsure{standart|standard} où le chang$^{\mathrm t}$ de base marche (mais où les foncteurs qui interviennent n'ont aucune raison d'être propres ou \unsure{lisses|lisse}).

Soient

\[ f:I\longrightarrow I',\qquad g:J\longrightarrow J' \]

de $\mathrm{Diag}$, d'où le diagramme (comparer page 33)

\diagram{\begin{tikzcd}[ampersand replacement=\&]
X'=I\times J \arrow[r,"G'"] \arrow[d,"{\deleted{\unsure{G}}F'}"'] \& I\times J'\qquad\unsure{X|X'} \arrow[d,"F"] \\
Y'=I'\times J \arrow[r,"G"] \& I'\times J'=Y
\end{tikzcd}}

Ce diagramme \deleted{\unsure{satis}} est tel que pour tt dér. à g. $\mathbb D$,
% archive sheet number: 52
